\begin{sectionbox}
	\subsection{Allgemein}
    \begin{tablebox}{ll}
        Molare Masse & $M$ \\
        Stoffmenge & $n=\frac{m}{M}$ \\
        Temperatur & $T[K]=T[^\circ C] + 273,15$ \\
        Freiheitsgrad & $f=3 \text{ (einatomiges Gas)}, =5 \text{ (zweiatomiges Gas)}$
    \end{tablebox}

	\begin{emphbox}
		\textbf{Ideale Gasgleichung:} $pV = nRT$ \\
        \textbf{Innere Energie:} $U=\frac{f}{2}nRT$
	\end{emphbox}

    \textbf{Wärmemenge Q im System} \\
	$Q=cm\Delta T$ \\ $c$: spez. Wärmekapazität \\

	\textbf{Gesetz von Boyle + Mariotte} \\
	für $\text{Teilchenanzahl}=\ir const.$ und $T=\ir const.$ gilt:\\
    $p_1V_1=p_2V_2$ \\

	\textbf{2. Gesetz von Gay Lussac} \\
    \begin{tabular}{l|l}
        \hline feste Gasmenge; $V=\ir const.$ & fester Druck; $p=\ir const.$ \\
        $\frac{p_1}{T_1}=\frac{p_2}{T_2}$ & $\frac{V_1}{T_1}=\frac{V_2}{T_2}$ \\ \hline
    \end{tabular}
\end{sectionbox}

\begin{sectionbox}
	\subsection{Hauptsätze der Thermodynamik}

	\begin{cookbox}{}
		\item[0.] Zwei Körper im thermischen Gleichgewicht zu einem dritten\\ $\rightarrow$ Alle stehen untereinander im Gleichgewicht
		\item[1.] $\Delta U = \Delta Q + \Delta W \rightarrow $ Es gibt kein Perpetuum mobile erster Art - Maschine mit $>$100\% Wirkungsgrad\\
		$\textbf{Verschiedene Möglichkeiten für Zustandsänderung:}$
		\begin{sectionbox}
			$\textbf{a)}$ Isobarer Prozess, $p =$ const.\\ $\rightarrow$ im idealen Gas ist $C_p$ konstant $\Rightarrow Q_{12} = C_p \Delta T$\\
			$\textbf{b)}$ Isochorer Prozess: $V =$ const.\\ $\rightarrow$ im idealen Gas ist $C_v$ konstant $\Rightarrow Q_{12} = \Delta U$\\
			$\textbf{c)}$ Isothermer Prozess: $T =$ const.\\ $\Rightarrow$ $W_{12} = -Q_{12}$
			$ =n RT \ln\frac{V_1}{V_2}$
			\\Freiwerdende Wärme: $Q_{12} = -W_{12}$\\
			$\textbf{d)}$ Adiabatischer Prozess: $\Delta Q = 0$\\
		\end{sectionbox}
		\item[2.] Thermische Energie ist nicht in beliebigem Maße in andere Energiearten umwandelbar. $\eta < 1$
	\end{cookbox}
\end{sectionbox}

TODO: Hauptsätze genauer, Adiabatengleichung, Entropie